% 国赛论文主文件（cumcmthesis 模板）—— 标准骨架
% 编译：xelatex ×2（先复制 cumcmthesis.cls 到本目录；图片复制到本目录或改相对路径）
% 来源：知识库/模板/国赛LaTeX/
\documentclass{cumcmthesis}

\title{论文题目：基于……的……研究}
\tihao{A}                    % 题号 A/B/C/D/E
\baominghao{参赛队号 12 位}
\schoolname{学校全称}
\membera{成员一}
\memberb{成员二}
\memberc{成员三}
\supervisor{指导教师}
\yearinput{2026}
\monthinput{09}
\dayinput{13}

\begin{document}

\maketitle

\begin{abstract}
在正文摘要处填写（800-1000 字，每问给"方法名 + 带单位数值"，无图表）。第一段总述问题+统一框架；其后每问一段，以"针对问题X"开头，含方法+核心数值；结尾总结升华（创新点+检验）。
\keywords{关键词一 \quad 关键词二 \quad 关键词三}
\end{abstract}

% 注意：2026 规范要求正文不要目录，切勿使用 \tableofcontents

\section{问题重述}
（原创语言重述：背景、核心目标、约束；提炼子问题，不复制原题）

\section{问题分析}
（"总—分"：先整体思路，再各问：关键矛盾→解题思路→核心步骤）

\section{模型假设}
\begin{enumerate}
    \item 假设一（合理性依据）……
    \item 假设二（合理性依据）……
    % 4-6 条，每条给依据，别过度简化
\end{enumerate}

\section{符号说明}
\begin{table}[H]
\centering
\caption{主要符号说明}
\begin{tabular}{ccc}
\toprule[0.8pt]
符号 & 含义 & 单位 \\
\midrule[0.5pt]
$D$ & 示例变量 & kg \\
\bottomrule[0.8pt]
\end{tabular}
\end{table}

\section{问题一的模型建立与求解}
\subsection{模型建立}
（公式用 equation 环境编号，变量先定义后使用）
\subsection{模型求解}
（写明数据/算法/软件，关键结果）
\subsection{结果分析}
（基础数值解读 + 深层机理 + 直接回应问题）
\subsection{模型检验}
（误差/灵敏度/稳健性，按问题需要）

\section{问题二的模型建立与求解}
\subsection{模型建立}
……
\subsection{模型求解}
……
\subsection{结果分析}
……

\section{问题三的模型建立与求解}
\subsection{模型建立}
……
\subsection{模型求解}
……
\subsection{结果分析}
……

\section{问题四的模型建立与求解（若有）}
……

\section{模型评价}
\subsection{模型优点}
（创新性/适用性/效率，结合数据）
\subsection{模型缺点与改进}
（假设限制/数据依赖/扩展局限，具体不空泛）

\section{参考文献}
\begin{thebibliography}{9}
\bibitem{1} 作者. 题目[J]. 期刊, 年份, 卷(期): 页码.
\end{thebibliography}

\appendix
\section{支撑材料说明}
（代码文件清单、中间结果、补充图表说明；2026 要求若用 AI 需附 AI工具使用详情.pdf）

\end{document}
